% Paste into Overleaf: Discussion chapter.
% Cross-refs assume Ch5 labels sec:results-rq1, sec:results-rq2, sec:results-errors, tab:results-b1.

\chapter{Discussion}
\label{ch:discussion}

The scenario contrasts in Chapter~\ref{ch:resultsAndAnalysis} show a split between final category match and intermediate control of belief and clarification.
This chapter interprets that split for RQ1 and RQ2, then turns to validity threats and concrete improvement directions.

\section{Intention conditioning}
\label{sec:discussion-intention}

Under RQ1, enabling the multi-turn intention path does not raise category task match relative to LLM-only operation (Section~\ref{sec:results-rq1}).
The useful signal sits in the process measures: posterior entropy falls more sharply, and clarification is triggered more often when the session remains underspecified.
In a booking setting where wrong early commitment is costly, that pattern is not empty: the sidecar makes uncertainty visible and gives the dialogue policy a reason to ask before suggesting.
At the same time, the higher clarify rate does not convert into better finals on the reported package.
Several mismatched sessions still contain clarify acts, which suggests that asking alone cannot repair thin or cross-class lexical evidence.
Intention conditioning should therefore be read as a control and inspectability layer rather than as an automatic accuracy upgrade over a language-model-only decision path.

The boundary errors summarised in Section~\ref{sec:results-errors} reinforce the same reading.
Wellness--fitness, beauty--medical, and dining or travel openings that share lifestyle vocabulary are exactly the places where a bag-of-words front-end and a category-dominant belief update are brittle.
A modular stack can expose those failures; it does not, by itself, remove the underlying class overlap.

\section{Memory and privacy}
\label{sec:discussion-memory}

Under RQ2, adding an aggregate-memory prior on the same intention path yields a modest match gain and a lower clarify rate, while the prompt privacy audit remains perfect (Section~\ref{sec:results-rq2}).
That combination matches the intended privacy claim more closely than it matches a strong personalisation claim.
Profiles and identifiers stay outside the generative prompt; memory influences generation only through the compact posterior summarised in the intention block.
The behavioural lift, however, depends on a mild gold-aligned category prior used for fair ablation.
It shows that a local prior can stabilise borderline sessions, not that long-horizon preference learning from real bookings has been demonstrated.

Nested reading remains important here.
Memory-augmented operation is Intention-conditioned operation plus a prior, not a separate stack whose score can be added on top of Intention-conditioned match.
The right comparison for RQ2 is therefore the presence or absence of that prior under an otherwise shared path, together with the privacy audit.

\section{Validity and reliability}
\label{sec:discussion-validity}

Several threats limit how far the reported contrasts should be generalised.
The dialogue package is scripted: openings, revisions, and topic changes are designed to probe vagueness, but they are not a sample of live Fixmeapp traffic.
Category task match is a surrogate for goal attainment in the absence of a booking API or database ground truth; Inform/Success-style completion scores are out of scope by construction.
The primary quantities describe the belief and dialogue-policy path; they are not scores of surface wording quality when a generative model is attached.

The Memory-augmented contrast is sensitive to how the prior is seeded.
A gold-aligned mild prior is appropriate for an ablation that asks whether memory can matter under privacy constraints, but it overstates what would be available from a cold start or from noisy historical aggregates.
Extractor errors on class boundaries propagate into belief and policy; scenario-level gains cannot outrun that front-end frontier indefinitely.
Reliability is supported by fixed scripts, frozen extractor artefacts, and deterministic offline replay; decoding variability would re-enter if generative replies were scored as a primary outcome, which they are not in the reported tables.

\section{Improvement directions}
\label{sec:discussion-improvements}

The results point to a small set of high-leverage improvements rather than to a general call for larger models.
Stronger separation among wellness, fitness, beauty, and medical cues would attack the dominant residual families more directly than further raising clarify frequency.
Clarification content could be tied more tightly to competing high-mass categories, so that ask acts collect the evidence the belief update actually needs.
Memory should move from gold-aligned seeds toward priors estimated from held-out interaction histories or from aggregate stores that are wrong as often as real profiles are, while keeping the same prompt allow-list.
Finally, any later study that scores user-visible replies should keep the sidecar metrics and the generative metrics separate, so that gains in fluency are not mistaken for gains in intention tracking.

